\begin{proof}
First consider the set of strings $\Xcal \subseteq \{0,1\}^{\tau n/10}$ with Hamming weight
\[
|x|
= \lfloor \tau n/100 \rceil \pm 1
< \tau n/20.
\]
By \cref{lem:adj-min-cut} the graph $G_x$, $x \in \Xcal$, has a unique minimum cut with shores $(V_1 \cup V_2, V_3 \cup V_4)$ and weight $2|x|$.
Now let $\Gcal' = \{G_x : x \in \Xcal\}$ and assume the existence of a quantum algorithm that for every $G_x \in \Gcal'$ computes with probability at least $2/3$ the weight $2|x|$ of a minimum cut in $G'$ with at most $q$ queries to the adjacency array of $G_x$.
By \cref{lem:adj-query} this is equivalent to outputting the Hamming weight $|x|$ with probability at least $2/3$ for any $x \in \Xcal$ while making only $q$ queries to $x$.
Using \cref{thm:thresh} this implies the lower bound $q \in \Omega(\tau n)$.

Next consider the set of strings $\Xcal' \subseteq \{0,1\}^{\tau n/10}$ that have Hamming weight $|x| = \tau n/20 \pm 1$.
By \cref{lem:adj-min-cut} the graph $G_x$, $x \in \Xcal'$, again has a unique minimum cut.
If $|x| = \tau n/20 - 1$ then its shores are $(V_1 \cup V_2, V_3 \cup V_4)$, while if $|x| = \tau n/20 + 1$ then its shores are $(V_1 \cup V_4, V_2 \cup V_3)$.
Now assume that there exists a quantum algorithm that with probability at least $2/3$ returns the shores of a minimum weight cut of $G_x$ with at most $q$ queries to the adjacency array of $G_x$.
By \cref{lem:adj-query} this is equivalent to distinguishing $|x| = \tau n/20 - 1$ from $|x| = \tau n/20 + 1$ with probability at least $2/3$ for any $x \in \Xcal$ while making only $q$ queries to $x$.
Using \cref{thm:thresh} this implies the lower bound $q \in \Omega(\tau n)$.
\end{proof}
